\title{LongRoPE: Extending LLM Context Window Beyond 2 Million Tokens}

\begin{abstract}
An extensive context window is a valuable attribute for large language models (LLMs). Nevertheless, present context window extensions rarely surpass approximately 128k tokens, primarily because of the substantial expense associated with fine-tuning, limited availability of sufficiently long training documents, and the emergence of problematic values caused by newly introduced token positions.

This work presents LongRoPE, a novel approach that, for the first time, scales the context window of pre-trained LLMs up to an exceptional \textbf{2048k} tokens, utilizing no more than 1k fine-tuning steps at sequence lengths up to 256k—while still preserving the model’s performance at its initial short context window. This advancement is accomplished through three central innovations: (i) we systematically uncover and leverage two types of non-uniformities inherent to positional interpolation via an efficient search process, resulting in superior fine-tuning initialization and allowing for an 8$\times$ context window increase even without additional fine-tuning; (ii) we propose a stepwise extension procedure, initially fine-tuning a model at 256k context length, followed by an additional stage of positional interpolation on this extended model to ultimately realize a 2048k context window; (iii) we perform a targeted adjustment of LongRoPE at 8k tokens to restore optimal short context window accuracy. Comprehensive experiments conducted on LLaMA2 and Mistral, spanning multiple tasks, verify the robustness and efficiency of our methodology. Notably, LLMs extended with LongRoPE maintain their original network structure with only minimal changes to their positional embedding component, and retain compatibility with the majority of existing optimization techniques. The codebase will be released at \texttt{https://github.com/microsoft/LongRoPE}

\section{Introduction}

Although Large Language Models (LLMs) have demonstrated impressive performance across a range of tasks~\cite{gpt4,llama2}, they are constrained by relatively short context window sizes; for instance, LLaMA2 can process only up to 4096 tokens~\cite{llama2}. When input extends beyond this window, model accuracy deteriorates because the LLM has not been trained to handle additional token positions. This limitation is particularly problematic in key applications such as in-context learning with many demonstrations~\cite{Huang2023FewerIM} and in the deployment of LLM agents~\cite{park2023generative,madaan2023selfrefine}.

Recent studies have demonstrated that it is possible to increase the context window of pre-trained LLMs up to approximately 128k by applying fine-tuning techniques on longer sequences~\cite{longlora,pi,yarn,zhang2024soaring,liu2023scaling}. Despite this progress, pushing the context window further encounters three principal challenges. To begin with, untrained positional indices corresponding to new positions introduce numerous catastrophic values, which in turn cause out-of-distribution complications and render the fine-tuning process difficult to converge~\cite{pi}. This problem is exacerbated when scaling from 4k to over 1000k context window sizes, as this transition adds more than 90\% novel positions. Additionally, effective fine-tuning typically demands access to texts matching the target length; unfortunately, currently available datasets provide few texts of such length, particularly those surpassing 1000k tokens. Beyond this, training with ultra-long texts substantially increases computational demands, necessitating vast training durations and large-scale GPU investments. Finally, stretching the context window to extreme lengths causes attention to be distributed across many more token positions, reducing its focus and ultimately diminishing the model’s performance on shorter contexts as originally designed~\cite{pi}.

A common strategy to address the initial challenge involves using interpolated RoPE positional embeddings~\cite{rope,pi}, where newly introduced position indices are scaled down to fit within the range seen during pretraining, as illustrated in Fig.~\ref{fig:overview}. Position Interpolation (PI)~\cite{pi} achieves this by linearly interpolating the rotary angles of RoPE in proportion to the extension factor. The NTK method~\cite{ntk,dynamicntk} takes a different approach, performing non-uniform interpolation and extrapolation along the various RoPE dimensions. YaRN~\cite{yarn} further subdivides the RoPE dimensions into three groups based on frequency characteristics, then applies a combination of extrapolation, NTK, and linear interpolation strategies to these groups. Nonetheless, the nature of positional embedding within the Transformer model exhibits \emph{intricately varying information entropy} that is not uniform across all positions. Present techniques do not fully exploit this non-uniformity, which results in partial loss of information and constrains the achievable context window length.

Section~\ref{sec:analysis} presents two main empirical observations: \textbf{(1)} Successful positional interpolation must account for two types of non-uniformities, specifically, differences across RoPE dimensions as well as variations in token positions. Minimal interpolation is advantageous for lower RoPE dimensions and for tokens near the beginning, although the most suitable strategies are influenced by the desired extended sequence length.
  \textbf{(2)} By integrating these non-uniform characteristics into the design of positional interpolation, it becomes possible to better preserve critical information inherent in the original RoPE—especially within essential dimensions and initial token slots. This approach significantly reduces information loss attributed to the interpolation process, resulting in superior initialization for subsequent fine-tuning. Additionally, such consideration supports up to an 8$\times$ length extension even in the absence of fine-tuning.

Inspired by these results, we propose \textbf{LongRoPE}, a robust approach capable of pushing the LLM context window to lengths surpassing 2 \emph{million} tokens. The design of LongRoPE revolves around three main innovations. To start, {LongRoPE} leverages the multidimensional and non-uniform characteristics present in positional interpolation. It determines suitable rescale factors for the rotation angles within RoPE for each separate dimension, guided by token position information. Given that the parameter space for these rescale factors grows exponentially in relation to the desired extension ratio, {LongRoPE} incorporates an evolutionary search strategy complemented by two optimization methods to improve the efficiency of this process. Figure~\ref{fig:overview} offers an illustration of the resultant rescaled RoPE discovered through this search.

Subsequently, {LongRoPE} utilizes a resource-efficient and incremental extension methodology to expand the context window up to 2048k, circumventing the necessity for direct fine-tuning on extremely lengthy text sequences, which are both rare and challenging to acquire. The process initiates by locating an optimal configuration for a 256k sequence length on the pre-trained LLM, followed by targeted fine-tuning at this specific length. Given that our non-uniform positional interpolation facilitates an 8$\times$ extension absent additional fine-tuning, a second phase involves identifying updated RoPE rescale parameters using the previously fine-tuned, lengthened LLM. Through this approach, we ultimately enable LLaMA2 and Mistral~\cite{mistral} models to support a 2048k context window.

To address potential drops in performance for the original, shorter context windows, {LongRoPE} retains flexibility in the adjustment of RoPE rescale factors in the extended LLM. Analogous to the upscaling from 256k to 2048k, we employ our search algorithm to downscale from the 256k fine-tuned model to shorter context lengths of 4k and 8k, thereby reducing the degree of positional interpolation required. At inference time, whenever the input sequence is under 8k tokens, RoPE is updated to incorporate the optimized rescale factors identified through our search.

A comprehensive series of experiments conducted on multiple LLMs and a range of long-context evaluation tasks highlights the efficacy of our proposed approach. Our results indicate that LongRoPE consistently preserves low perplexity across evaluation lengths spanning from 4k to 2048k, achieves passkey retrieval accuracy exceeding 90\%, and matches the accuracy of existing models on established benchmarks with a 4096 token context window. The LongRoPE technique is compatible with all LLMs utilizing RoPE embeddings. We intend to make both our codebase and the LongRoPE-2048k models publicly available.

\section{Non-uniformity in Positional Interpolation}
\label{sec:analysis}

\subsection{Preliminary}

Transformer models require explicit positional information, often in the form of position embedding, to represent the order of input tokens. Our work focuses on the RoPE~\cite{rope} position embedding, which is widely used in recent  LLMs. For a token at position index $n$, its corresponding RoPE encoding can be simplified as follows:

\begin{equation}
		\fontsize{7.8}{7.8} \selectfont
	\label{eq:rope}
	\left[ cos(n\theta_0), sin(n\theta_0), cos(n\theta_1), \cdot\cdot\cdot, cos(n\theta_{d/2-1}), sin(n\theta_{d/2-1}) \right]   
\end{equation}

where $d$ refers to the dimension of the embedding, 
$n\theta_i$ indicates the rotary angle associated with a token at the $n$-th position,
and $\theta_i=\theta^{-2i/d}$ specifies the frequency of rotation. The RoPE method typically utilizes a default $\theta$ base value of 10000.

\textbf{\emph{Context window extension ratio $s$} and positional interpolation}. The parameter $s$ is introduced as the proportion between the expanded context length $L'$ and the initial context length $L$, calculated as $s=\frac{L'}{L}$.

To extend the context window from $L$ to $L'$, current positional interpolation methods suggest downscaling rotation frequencies $\theta_i$ based on extension ratio $s$. Let $\beta=\theta^{2/d}$, and $\lambda$ denote the actual rescale factor related to $s$, we   unify these positional interpolation methods as follows:

\begin{equation}	
	\fontsize{7.5}{7.5} \selectfont
	\label{eq:unified_rope}
	\begin{aligned}
		\left[cos\left(\frac{n}{\lambda(\beta)^0}\right), sin \left(\frac{n}{\lambda(\beta)^0}\right), cos\left(\frac{n}{\lambda(\beta)^1}\right), \cdot\cdot\cdot,  sin \left(\frac{n}{\lambda(\beta)^{d/2-1}}\right) \right]   
	\end{aligned}
\end{equation}

\textbf{Linear positional interpolation (PI)}. The PI approach~\cite{pi} proposes to interpolate position indices linearly when extending sequences, as long as they remain within the maximum pre-trained length. When the sequence is extended by a ratio $s$, each RoPE rotation angle is uniformly scaled down by a factor of $\lambda=s$. This method causes positional encodings to become densely packed, which impairs the model's capability to differentiate between adjacent token positions. Consequently, PI generally exhibits suboptimal performance when dealing with large extension ratios.

\textbf{NTK-based interpolation and extrapolation}. In examining RoPE through the lens of information encoding, ~\cite{ntk,dynamicntk} leverage Neural Tangent Kernel (NTK) theory~\cite{jacot2018neural,tancik2020fourier}. To address the problem of position crowding inherent in PI, their approach involves distributing the interpolation demands over the various dimensions of RoPE. This method involves applying less scaling to the lower (high frequency) dimensions and greater scaling to the higher (low frequency) dimensions, facilitating both interpolation and extrapolation of positions, characterized by $\lambda=s^{i}$. The dynamic NTK model introduced in~\cite{dynamicntk} further enhances this by varying the extension ratio per position, adapting to the current sequence length. In contrast to PI, which typically requires fine-tuning, NTK-based approaches can achieve context window extension in scenarios that do not involve fine-tuning, although such methods commonly allow for a maximum extension up to 4$\times$.

\textbf{YaRN}~\cite{yarn} presents a notable advancement in the efficacy of positional interpolation methods. The approach partitions RoPE dimensions into three distinct categories based on frequency, and applies separate interpolation techniques to each. Specifically, dimensions associated with high frequencies are subject to extrapolation ($\lambda$=1), whereas those corresponding to low frequencies adopt linear interpolation (PI). Intermediate frequency dimensions are managed through the NTK approach. A central aspect of YaRN is its dimension grouping procedure, which is currently determined through manual empirical analysis. As a consequence, this reliance on human intervention can lead to less-than-optimal outcomes when applied to novel LLMs.

\subsection{Study on Non-uniform Positional Interpolation}

Drawing motivation from NTK and YaRN, we observe that their improvements stem from introducing non-linear components, notably through handling various frequencies along RoPE dimensions for tailored interpolation and extrapolation. Nonetheless, existing approaches to non-linearity predominantly depend on heuristics devised by humans. This observation leads us to pose two important questions: (1) Does the existing method for positional interpolation represent an optimal solution? (2) Might there be other, yet undiscovered, forms of non-linearity?

In order to address these questions, we employ evolution search (refer to Sec.~\ref{sec:method}) to identify improved non-uniform positional interpolation strategies for LLaMA2-7B. This search process is directed by perplexity measurements, utilizing 5 randomly selected samples from the PG19~\cite{pg19} validation set. Our experimental investigation yields several significant observations.

\textbf{Finding 1}: \textit{RoPE dimensions exhibit substantial non-uniformities, which are not effectively handled by current positional interpolation methods.}

We conduct a search for the optimal $\lambda$ value for each RoPE dimension as described in Eq.~\ref{eq:unified_rope}. Table~\ref{tbl:exp1} presents a comparison of LLaMA2-7B perplexity using various methods on the PG19 and Proof-pile~\cite{proof-pile} test datasets, all evaluated without fine-tuning. Our optimized configuration leads to clear gains in performance, indicating that existing linear (PI) and non-uniform (Dynamic-NTK, YaRN) interpolation strategies do not achieve optimal results. Interestingly, YaRN performs worse than both PI and NTK on PG19; this is attributable to its inability to utilize the full target context window in settings without fine-tuning. To illustrate, YaRN’s perplexity markedly increases beyond 7k tokens within an 8k context window.

Our investigation reveals that the rescaled coefficients $\lambda$ in Eq.~\ref{eq:unified_rope} are not uniform across dimensions, unlike the constant scaling factor $s$ utilized in PI, the formulaic approach of NTK, or the group-wise adjustment of YaRN. Employing these non-uniform scaling parameters leads to a notable enhancement in LLaMA2’s language modeling capabilities (specifically, in terms of perplexity) at 8k and 16k context lengths, all without the need for further fine-tuning. This improvement arises because the derived positional embeddings more faithfully maintain the characteristics of the original RoPE—most importantly, in essential dimensions—thereby easing the model's challenge in differentiating between positions of nearby tokens.

\textbf{Finding 2}: \textit{RoPE for the initial tokens in the input sequence should be extrapolated with less interpolation.} 

We propose that for the first $\hat{n}$ tokens in the input sequences, RoPE should perform minimal or no interpolation. This suggestion is based on evidence that these initial tokens typically acquire high attention scores and, consequently, play an essential role in the functionality of the attention layers, as reported in both Streaming LLM~\cite{streamingllm} and LM-Infinite~\cite{han2023lminfinite}. To test this idea, we increase the context window size to 8k and 16k using PI and NTK approaches, selectively omitting interpolation for the first $\hat n$ tokens (where $\hat n$ varies as 0,2, ..., 256). Setting $\hat n$=0 results in the default PI and NTK application. Table~\ref{tbl:tokenposition} presents two main findings: \textbf{(1)} excluding position interpolation for the initial tokens leads to enhanced performance, and \textbf{(2)} the ideal value of $\hat n$—the number of initial tokens to exclude—varies in relation to the intended extension length.

\textbf{Finding 3}: \textit{Non-uniform positional interpolation effectively extends LLM context window in both fine-tuning and non-fine-tuning settings.} 

Our experiments demonstrate that the application of our optimized non-uniform position interpolation leads to substantial enhancements in extension capabilities at 8k and 16k context lengths even without additional fine-tuning. However, pushing these extensions further necessitates fine-tuning. Therefore, we further train LLaMA2-7B equipped with our optimized RoPE for a 64k context window (refer to the Appendix for experimental details). As indicated in Table~\ref{tbl:finetune}, our approach surpasses both PI and YaRN considerably, both prior to and following fine-tuning the LLaMA2-7B model. This performance boost can be attributed to the judicious exploitation of non-uniform positional interpolation, which curtails information loss and results in a superior starting point for fine-tuning.

\textbf{Summary}. This work identifies two primary sources of non-uniformity: the variation among RoPE dimensions and across token positions. By harnessing these non-uniformities within positional interpolation, our approach substantially elevates context extension outcomes in large language models.

\section{LongRoPE}

\label{sec:method}
Building on these insights, we propose LongRoPE—a method that begins by implementing an effective search algorithm designed to leverage both forms of non-uniformity. Utilizing this algorithm, LongRoPE successfully expands the LLM context window to accommodate over 2 million tokens.

\subsection{Problem Formulation}

These dual non-uniformities expand the solution space considerably and add layers of complexity to the optimization process. To manage these challenges, we recast the multidimensional non-uniform position interpolation optimization as a search problem.

For a LLM targeting a  context window size of $L'$ and lengthy input documents $\mathbf{X}$, where each $\mathbf{x}\in\mathbf{X}$ surpasses $L'$ in token length, we denote the original rotary angle of the $i^{th}$ dimension in RoPE embedding at token position $n$ as  $\frac{n}{\beta^i}$. The optimization problem is then formulated as follows:

\begin{equation}
\begin{gathered}
		\label{eq:problem}

\underset {\mathbf{x}\in\mathbf{X}; \, |\mathbf{x}|\ge L'}{ \operatorname {arg\,min} } \,  \mathcal{L}\left( \text{LLM} (\text{RoPE}, \mathbf{X})\right), \quad \text{where}
\\
\scalebox{0.93}{
$\underset {\begin{subarray}{l} i=0,\cdots,\frac{d}{2}-1;\\ n\in[ 0,|\mathbf{x}|);\end{subarray}}{\text{RoPE}_(n)}= \left[\cdots, \cos\left(\mathbb{I}(\hat{\lambda}_{i}, \hat{n})\frac{n}{\beta^i}\right), \sin\left(\mathbb{I}(\hat{\lambda}_{i}, \hat{n})\frac{n}{\beta^i}\right), \cdots\right]$}
\\
\text{with} \; \mathbb{I}(\hat{\lambda}_{i},\hat{n})=
\begin{cases}
1 &\text{if } n<\hat{n} \\
\frac{1}{\lambda_{i}} &\text{if } n \geq \hat{n}
\end{cases}
\end{gathered}
\end{equation}

In this context, we define a collection of scaling coefficients, $\mathbb{I}(\hat{\lambda}_{i},\hat{n})$, that address both types of non-uniformities present. Here, $\hat{\lambda}_{i}$ represents the adjustment for each RoPE dimension’s non-uniformity, while $\hat{n}$ indicates the threshold token position at which the modification is introduced. The function $\mathbb{I}(\hat{\lambda}_{i},\hat{n})$ is utilized to adjust the rotational angle corresponding to the $i^{th}$ RoPE dimension. Notably, for token positions up to $\hat{n}-1$, the scaling parameter $\hat\lambda_{i}$ does not influence the computation, and the standard RoPE rotation angle $\frac{n}{\beta^i}$ is retained. Once token positions satisfy $n\ge\hat{n}$, however, the rescale factor is enacted.

Assuming a desired context window length of $L'$, our goal is to determine the best rescale factors, represented as ($\mathbb{I}(\hat{\lambda}_{0},\hat{n})$, $\mathbb{I}(\hat{\lambda}_{1},\hat{n})$, ..., $\mathbb{I}(\hat{\lambda}_{i},\hat{n})$, ...), for each RoPE dimension from the first to the $d$th. This adjustment should allow the target $\text{LLM}$, equipped with the modified $\text{RoPE}$, to minimize the next-token prediction loss $\mathcal{L}$ (which corresponds to perplexity) across input sequences $\mathbf{X}$ of length $L'$.

\subsection{Searching the Non-uniform Position Interpolation}

In addressing the challenge described in Eq.~\ref{eq:problem}, we present a straightforward but powerful approach that identifies optimal RoPE rescale factors, thereby leveraging the complete range of multidimensional non-uniformities inherent in position embedding.

\textbf{Search space}. To account for the two types of non-uniformity, we construct an extensive search space. Table~\ref{tbl:searchspace} provides details regarding this search space. In particular, our approach permits tuning a unique rescale factor for each dimension within the RoPE embedding. For ease in structuring the search space, we opt to explore values for $\lambda_i$ and $\hat{n}$ rather than directly searching over $\mathbb{I}(\hat{\lambda}_{i},\hat{n})$, with the relationship $\hat{\lambda}_{i}=1/\lambda_i$. As presented in Table~\ref{tbl:searchspace}, $\lambda_i$ is explored within a range starting at 1.0 (which corresponds to straightforward extrapolation) and extending up to $s\times1.25$ (which allows for broader interpolation compared to PI), in increments of 0.01, where $s$ denotes the desired extension ratio of the context window.

The parameter $\hat{n}$ determines how many leading token positions maintain their original RoPE embedding, without applying position interpolation. In practice, we explore $\hat{n}$ values from the set \{0, 1, 2, 4, 8, 12, 16, 20, 24, 28, 32, 64, 128, 256\}. If $\hat{n}=0$, rescale factors are applied to every token position.

\textbf{Evolution-based search}. The search space described in Table~\ref{tbl:searchspace} encompasses a wide array of positional interpolation strategies, making systematic exploration computationally demanding. As an illustration, applying a $s=4\times$ extension results in $400^{128/2}\times14 = 4\times10^{167}$ possible configurations. This search space grows exponentially with increased extension ratios. To tackle this complexity, we adopt evolution search~\cite{guo2020single} and incorporate two optimization strategies that significantly enhance the efficiency of the search process. The general flow of our search algorithm is provided in Algorithm~\ref{alg:search}.

\textit{Optimized initial population generation}. Rather than relying solely on random initialization for the population of $P$ rescale factors, we explicitly incorporate the three RoPE rescale factors associated with PI, NTK, and YaRN into the initial group as individual candidates. The other $P$-3 members of the population are created by applying random mutations to these three rescale factors, each with probability $p$.

\textit{Monotonically non-decreasing constraint}. Following the generation of the initial population, we assess each individual by measuring LLM perplexity. For this, the relevant RoPE rescale factors are applied within the target LLM, and perplexity is evaluated for the input $\mathbf{X}$. The individuals with the lowest perplexity (top-k) are selected to serve as parents for subsequent evolutionary steps. Nonetheless, because the search space is highly extensive, straightforward approaches to mutation and crossover may result in suboptimal candidates, causing unnecessary perplexity evaluations. This inefficiency becomes especially pronounced when $L'$ is large, as each perplexity estimation entails substantial inference overhead.

To tackle this, we enforce a monotonic non-decreasing constraint on the sequence of sampled RoPE rescaling factors such that $\lambda_i \le \lambda_{i+1}$. We exclusively use RoPE parameterizations meeting this requirement during perplexity assessments on LLMs, which substantially curtails the search space. Concretely, we stipulate that $\lambda_i$ must monotonically increase as the RoPE dimension index $i$ progresses from $0$ through $63$. This dimension-based monotonicity is motivated by NTK theory~\cite{jacot2018neural,tancik2020fourier,ntk}, which posits that lower-indexed dimensions, associated with higher frequencies, benefit from reduced interpolation (i.e., smaller values of $\lambda_i$), while higher-indexed dimensions, corresponding to lower frequencies, are more suitable for increased interpolation (i.e., larger $\lambda_i$).

\begin{algorithm}[t]
	\small
	\caption{The search algorithm}
	\textbf{Input:} target LLM, input samples $\mathbf{X}$, population size $P$, mutation size $N_1$, crossover size $N_2$, max iterations $\mathcal{T}$, mutate probability $p$\\

	\begin{algorithmic}[1]
		\label{alg:search}
		\STATE Top-k$\leftarrow \varnothing$;	
		\STATE $\text{P}_0$ := \textit{Initialize\_population\_with\_optimization}($P$, $\mathbf{X}$, $p$); 
		\FOR{$i=1$ to $\mathcal{T}$}
		\STATE \textit{Compute\_perplexity}(LLM, $\text{P}_{i-1}$, $\mathbf{X}$);
		\STATE Top-k $\gets$ \textit{Update\_Topk}(Top-k, $\text{P}_{i-1}$);
		\STATE $\text{P}_{mutation}$ := \textit{Mutation\_with\_mono\_constraint}(Top-k, $N_1$, $p$); 
		\STATE $\text{P}_{crossover}$ := \textit{Crossover\_with\_mono\_constraint}(Top-k, $N_2$);
		\STATE $\text{P}_i$ := $\text{P}_{mutation} \cup \text{P}_{crossover} \cup$ Top-k;
		\ENDFOR
		\STATE Output the candidate in Top-k that yields the minimum perplexity;
	\end{algorithmic}
\end{algorithm}

\textbf{8$\times$ extension without fine-tuning}. Through evolutionary search, our approach successfully discovers non-uniform RoPE rescale factors that maintain crucial dimensions and sequence positions, thereby reducing information degradation from interpolation. As shown in Fig.\ref{fig:non-ft}, this technique allows LLaMA2's context window to scale from 4k up to 32k tokens with no need for additional fine-tuning. By comparison, alternative methods like PI, as well as non-uniform NTK and YaRN, experience a pronounced increase in perplexity once the context window is doubled.

\subsection{Extending LLM Context Window to 2048K}
\label{sec:extendto2048k}

\textbf{Progressive extension to 2048k}. In this section, we present our strategy for expanding the context window of pre-trained LLMs from the standard 4k up to beyond 2048k. Our approach leverages non-uniform positional interpolation, which has been shown to enable up to an 8$\times$ increase in context length without the need for fine-tuning. When even greater expansion is desired (such as a 512$\times$ increase), fine-tuning becomes necessary. A possible approach involves determining appropriate RoPE rescaling factors for the target 2048k window, followed by a fine-tuning stage. Nonetheless, this process is hindered by the significant computational demands required for training. Furthermore, our observations indicate that effective fine-tuning for such large extension ratios is particularly challenging (refer to Appendix for details).

Fortunately, LongRoPE demonstrates efficacy when applied to both base and fine-tuned LLMs with extended context. To this end, we propose a progressive and efficient approach capable of reaching the intended 2048k target using merely 1k fine-tuning steps, while confining the training context length to 256k.

$\Diamond$ \textit{Searching for 256k pre-trained LLM extension using LongRoPE}. For instance, using LLaMA2, we perform a parameter search targeting context window sizes of 128k and 256k. Here, the scaling factors used for the extension process are 32$\times$ and 64$\times$, accordingly.

$\Diamond$ \textit{Fine-tuning to 256k}. Next, we adapt the pre-trained LLM to reach a context window of 256k through a staged fine-tuning approach. Initially, we fine-tune LLaMA2 for 400 steps utilizing RoPE rescaled factors set for 128k. Afterward, we switch to RoPE rescaled factors for 256k on the resulting checkpoint and proceed with an extra 600 steps of fine-tuning. This sequential fine-tuning strategy demonstrates greater efficiency compared to attempting direct fine-tuning at 256k.

$\Diamond$ \textit{Scaling fine-tuned extended LLM up to 2048k via LongRoPE search}. In the final step, we apply a secondary search process on the LLM previously fine-tuned for 256k sequence lengths. This method expands the context window to an exceptional size of 2048k tokens, all achieved without any additional fine-tuning. Consequently, the context length is extended by a factor of $512\times$.

\textbf{Performance within original context window}. Upon increasing the context window length to an extensive 2048k, we observe a degradation in performance for sequences confined to the initial context window. This phenomenon, previously documented in the context of positional interpolation~\cite{pi}, arises because mapping positional embeddings into higher-dimensional space causes the embeddings corresponding to positions in the original context window to cluster within a restricted area. This clustering diminishes the model’s capability within its typical context range. Especially when utilizing a 512$\times$ expansion factor, embeddings for positions in the initial 4k window are densely packed, exacerbating this effect.

To address this issue, we conduct an additional evolution search on the extended LLM, fine-tuning the RoPE rescale factors specifically for shorter context windows such as 4k and 8k. The upper limit for the searched $\lambda$ values is decreased, reflecting the reduced need for positional interpolation at these lengths. At inference time, the LLM automatically modifies the associated RoPE rescale factors in accordance with the context length.

\section{LongRoPE}

\label{sec:method}
Building on these insights, we propose LongRoPE, a method that initially incorporates a computationally efficient search algorithm designed to take full advantage of the identified non-uniformities. This approach is subsequently applied to extend the context window of LLMs to exceed 2 million tokens.

\subsection{Problem Formulation}

These two forms of non-uniformity result in an expansive array of possible solutions and add significant complexity to the optimization process. To tackle this challenge, we recast the optimization problem of multidimensional non-uniform position interpolation as a search problem.

For a LLM targeting a  context window size of $L'$ and lengthy input documents $\mathbf{X}$, where each $\mathbf{x}\in\mathbf{X}$ surpasses $L'$ in token length, we denote the original rotary angle of the $i^{th}$ dimension in RoPE embedding at token position $n$ as  $\frac{n}{\beta^i}$. The optimization problem is then formulated as follows:

\begin{equation}
\begin{gathered}
		\label{eq:problem}

\underset {\mathbf{x}\in\mathbf{X}; \, |\mathbf{x}|\ge L'}{ \operatorname {arg\,min} } \,  \mathcal{L}\, \left( \text{LLM} (\text{RoPE}, \mathbf{X})\right), \ \text{with} 
\\
\scalebox{0.93}{
$\underset {\begin{subarray}{l} i=0,\dots,\frac{d}{2}-1;\\ n\in[ 0,|\mathbf{x}|);\end{subarray}}{\text{RoPE}_(n)}= {\left[\ldots, \cos\left(\mathbb{I}(\hat{\lambda}_{i}, \hat{n})\cdot\frac{n}{\beta^i}\right),\ \sin\left(\mathbb{I}(\hat{\lambda}_{i}, \hat{n})\cdot\frac{n}{\beta^i}\right), \ldots\right] }$}\\
	\text{with} \quad \mathbb{I}(\hat{\lambda}_{i},\hat{n})=
\begin{cases}
	1&\mbox{\quad if $n<\hat{n}$}\\
	{\frac{1}{\lambda}_{i}}&\mbox{\quad if $n\geq\hat{n}$}
\end{cases}
	\end{gathered}
\end{equation}

In this approach, we define a group of scaling factors, $\mathbb{I}(\hat{\lambda}_{i},\hat{n})$, that address the two types of non-uniformities present. Here, $\hat{\lambda_i}$ pertains to non-uniformity among the RoPE dimensions, while $\hat{n}$ refers to variations across token positions. The function $\mathbb{I}(\hat{\lambda}_{i},\hat{n})$ is employed to adjust the rotation angle corresponding to the $i^{th}$ RoPE dimension, with $\hat{\lambda}_{i}$ serving as the scaling parameter and $\hat{n}$ indicating a token position threshold. For the first $\hat{n}-1$ token positions, no scaling is introduced; thus, the conventional RoPE rotary angle $\frac{n}{\beta^i}$ is maintained. Once token positions reach $n\geq \hat{n}$, the scaling factor becomes active and modifies the rotation accordingly.

Assuming a desired context window of length $L'$, the task is to determine the set of optimal rescale factors ($\mathbb{I}(\hat{\lambda}_{0},\hat{n})$, $\mathbb{I}(\hat{\lambda}_{1},\hat{n})$, ..., $\mathbb{I}(\hat{\lambda}_{i},\hat{n})$, ...) across the $1^{st}$ to $d^{th}$ dimensions of the RoPE module. Applying these rescaling adjustments to the target $\text{LLM}$ aims to minimize the next-token prediction loss, $\mathcal{L}$ (i.e., perplexity), on input data $\mathbf{X}$ of token length $L'$.

\subsection{Searching the Non-uniform Position Interpolation}

To address the issue presented in Eq.~\ref{eq:problem}, we propose a straightforward but powerful approach that identifies optimal RoPE rescale factors, thereby taking full advantage of the multidimensional disparities inherent in position embeddings.

\textbf{Search space}. We construct an expansive search space that accommodates both forms of non-uniformity. Table~\ref{tbl:searchspace} summarizes the details of this search space. In particular, our approach enables the search for unique rescaling factors corresponding to each individual dimension in the RoPE embedding. For the sake of simplicity, the search is conducted over the variables $\lambda_i$ and $\hat{n}$, rather than directly optimizing for $\mathbb{I}(\hat{\lambda}_{i},\hat{n})$, with $\hat{\lambda}_{i}=1/\lambda_i$. As indicated in Table~\ref{tbl:searchspace}, the range for $\lambda_i$ begins at 1.0 (which denotes pure extrapolation) and extends up to $s\times1.25$ (allowing for interpolation exceeding that of PI), incremented by 0.01. Here, $s$ denotes the extension ratio of the intended context window.

$\hat{n}$ determines how many of the initial token positions are kept with their original RoPE embeddings, foregoing position interpolation. In our experiments, we vary $\hat{n}$ over the set \{0, 1, 2, 4, 8, 12, 16, 20, 24, 28, 32, 64, 128, 256\}. If $\hat{n}=0$, then every token position adopts the rescale factors discovered during the search process.

\textbf{Evolution-based search}. The extensive search space described in Table~\ref{tbl:searchspace} encompasses a wide array of positional interpolation options, making efficient navigation highly nontrivial. As an illustration, a $s=4\times$ extension results in $400^{128/2}\times14 = 4\times10^{167}$ possible configurations. The combinatorial complexity increases exponentially with larger extension ratios. To contend with this, we employ evolution search~\cite{guo2020single} along with two additional optimization strategies to substantially improve search efficiency. The general procedure is outlined in Algorithm~\ref{alg:search}.

\textit{Optimized initial population generation}. Rather than relying solely on random initialization for the population of $P$ rescale factors, we explicitly include the three RoPE rescale factors associated with PI, NTK, and YaRN as distinct individuals at the start. The other $P$-3 individuals are produced by applying random mutations to these three rescale factors, with each mutation occurring with probability $p$.

\textit{Monotonically non-decreasing constraint}. Upon creating the initial population, we evaluate the LLM perplexity associated with each candidate. This involves adjusting the RoPE rescale factors within the target LLM, then measuring the perplexity on input $\mathbf{X}$. The $k$ best-performing candidates are selected as the parent set for the evolutionary process. Nonetheless, the expansive search domain means that unsophisticated mutation and crossover steps can frequently yield suboptimal candidates, resulting in superfluous perplexity evaluations. This inefficiency is further exacerbated for large $L'$, as computing each perplexity score requires substantial inference time.

To mitigate this issue, we enforce a non-decreasing monotonicity restriction on the sampled RoPE rescaling factors: $\lambda_i\le \lambda_{i+1}$. Only those RoPE configurations that fulfill this monotonicity criterion are utilized for evaluating LLM perplexity, thus substantially narrowing down the search space. In particular, this requirement means that $\lambda_i$ must strictly increase with respect to the RoPE dimension index ($i=0,\ldots,63$). This structural monotonicity aligns with insights from NTK theory~\cite{jacot2018neural,tancik2020fourier,ntk}, which posits that dimensions associated with higher frequencies (lower indices) necessitate less interpolation (smaller values of $\lambda_i$), while those corresponding to lower frequencies (higher indices) can accommodate increased interpolation (larger $\lambda_i$ values).

\begin{algorithm}[t]
	\small
	\caption{The search algorithm}
	\textbf{Input:} target LLM, input samples $\mathbf{X}$, population size $P$, mutation size $N_1$, crossover size $N_2$, max iterations $\mathcal{T}$, mutate probability $p$\\

	\begin{algorithmic}[1]
		\label{alg:search}
		\STATE Initialize Top-k $\leftarrow \phi$;
		\STATE $\text{P}_0$ $\leftarrow$ \textit{Initialize\_population\_with\_optimization}($P$, $\mathbf{X}$, $p$);
		\FOR{$i$ = $1$ \TO $\mathcal{T}$}
		\STATE \textit{Compute\_perplexity} (LLM, $\text{P}_{i-1}$, $\mathbf{X}$);
		\STATE Top-k $\leftarrow$ \textit{Update\_Topk}(Top-k, $\text{P}_{i-1}$);
		\STATE $\text{P}_{mutation}$ $\leftarrow$ \textit{Mutation\_with\_mono\_constraint}(Top-k, $N_1$, $p$);
		\STATE $\text{P}_{crossover}$ $\leftarrow$ \textit{Crossover\_with\_mono\_constraint}(Top-k, $N_2$);
		\STATE $\text{P}_i$ $\leftarrow$ $\text{P}_{mutation}$ $\cup$ $\text{P}_{crossover}$ $\cup$ Top-k;
		\ENDFOR
		\STATE Output the member from Top-k exhibiting the minimum perplexity;
	\end{algorithmic}
\end{algorithm}

\textbf{8$\times$ extension without fine-tuning}. Through evolutionary search, our approach successfully discovers non-uniform RoPE rescale factors that retain crucial dimensions and positional encodings, thereby reducing information loss caused by interpolation. As illustrated in Fig.\ref{fig:non-ft}, this technique enables LLaMA2 to scale its context window from 4k up to 32k tokens without necessitating fine-tuning. By comparison, alternative methods—including PI, as well as non-uniform NTK and YaRN—lead to a sharp increase in perplexity following a 2$\times$ context extension.

\subsection{Extending LLM Context Window to 2048K}
\label{sec:extendto2048k}

\textbf{Progressive extension to 2048k}. In this section, we present our approach for increasing the context window of pre-trained LLMs beyond the standard 4k limit, reaching up to over 2048k tokens. As shown previously, employing our non-uniform positional interpolation method enables us to extend the window by up to 8$\times$ without the need for fine-tuning. However, when attempting much larger extensions (e.g., by a factor of 512$\times$), fine-tuning becomes essential. One possible strategy involves optimizing the RoPE rescaling parameters for the desired 2048k context length and then performing fine-tuning. Nevertheless, this approach encounters significant difficulties, primarily due to the extremely high computational costs involved. Additionally, from our practical experience, effectively fine-tuning LLMs with such a substantial extension ratio presents considerable challenges (refer to Appendix).

Luckily, LongRoPE demonstrates efficacy with both base and fine-tuned, length-extended LLMs. As a result, we propose a streamlined, stepwise strategy that reaches the desired 2048k context window using only 1,000 fine-tuning iterations, with training restricted to a maximum length of 256k.

$\Diamond$ \textit{LongRoPE-based expansion of pre-trained LLM to 256k via search}. Using LLaMA2 as a case study, we perform a search process targeting context window lengths of 128k and 256k. Here, the context window is enlarged by a factor of 32$\times$ and 64$\times$, accordingly.

$\Diamond$ \textit{Fine-tuning to 256k}. Following pre-training, we further adjust the LLM to handle a 256k context window. The process involves initially fine-tuning LLaMA2 for 400 steps utilizing the RoPE rescaled factors tailored to 128k. Subsequently, the RoPE rescaled factors are switched to 256k at the obtained checkpoint, after which another 600 fine-tuning steps are performed. This staged approach is found to be more efficient compared to fine-tuning directly at the 256k window size.

$\Diamond$ \textit{Scaling up fine-tuned extended LLM to 2048k via LongRoPE search}. In the last step, we conduct an additional search process on the LLM that was previously fine-tuned for a 256k context length. Through this method, we achieve an exceptionally large context window of 2048k, circumventing any need for extra fine-tuning. This corresponds to a total extension factor of $512\times$.

\textbf{Recovery within shorter context windows}. Upon expanding the context window to a substantial length of 2048k, we observe diminished performance in the initial context window range. This phenomenon has been previously identified with positional interpolation~\cite{pi}, which compresses the position embeddings of the original window into a limited subset of the higher-dimensional space, thereby impairing the language model’s efficacy. When the extension ratio is increased to 512$\times$, the positional representations within the original 4k window become especially concentrated, exacerbating this issue.

To address this issue, we conduct an additional evolution-based search on the expanded LLM, specifically tuning RoPE rescale factors for shorter context lengths such as 4k and 8k. The upper bound for the searched $\lambda$ is lowered, since these shorter sequences demand less positional interpolation. At inference time, the LLM autonomously modifies the relevant RoPE rescale factors as needed.

 \section{Experiments}

\label{sec:exp}
\subsection{Setup}
\textbf{Evaluation Tasks and models}. LongRoPE is integrated with LLaMA2-7B and Mistral-7B models, whose effectiveness is assessed along three dimensions: (1) perplexity measured on long-form text to gauge the handling of extended contexts; (2) the Passkey retrieval task, which tests the model’s proficiency at extracting a designated passkey embedded amidst predominantly extraneous text; and (3) conventional LLM benchmark tasks constrained to the shorter 4096-token context window.

\textbf{Fine-tuning}. In the case of LLaMA2, a learning rate of 2e-5 is employed with a linear decay schedule and a total batch size of 32. The model is first fine-tuned for 400 steps using the Redpajama~\cite{redpajama} dataset, which is segmented into 128k-length pieces, each wrapped with BOS and EOS tokens. Subsequently, fine-tuning is continued from the resulting checkpoint for an additional 600 steps, extending to a 256k context window. Training with the 128k context involves 8 A100 GPUs utilizing the distributed training framework~\cite{cube}; extending to 256k requires scaling to 16 A100 GPUs. For Mistral, a fixed learning rate of 1e-6 and a total batch size of 64 are set. Both the 128k and 256k context window models are fine-tuned according to the YaRN configuration~\cite{yarn}, performing 400 training steps on Together Computer’s Long-Data Collections~\cite{mistraldata} at a sequence length of 16k. This process leverages 4 A100 GPUs.

\textbf{Search}. For target window sizes up to 256k, the parameters are set as follows: $P$=64, $N_1$=$N_2$=16, $p$=0.3, and $\mathcal{T}$=40. In each iteration, the top 32 candidates are chosen for mutation and crossover. Perplexity is assessed using 5 randomly selected validation samples from PG19, ensuring that each sample meets or exceeds the specified context length. When the window size exceeds 512k, we reduce by half the sizes for population, mutation, and crossover. In this case, perplexity evaluation is conducted using 3 randomly drawn samples from the Pile-Books3~\cite{pile} validation dataset.

\textbf{Baselines}. In order to achieve a 2048k context length, we performed fine-tuning on models using context windows of 128k and 256k, resulting in LongRoPE-2048k (ft=128k) for LLaMA2 and LongRoPE-2048k (ft=256k) for Mistral. We benchmark these four models against leading context window extension baselines. In particular, we focus on publicly available LLMs that have been fine-tuned following positional interpolation approaches such as PI, NTK, and YaRN. The comparison set includes Together-32k~\cite{together}, Code LLaMA~\cite{codellama}, LongLoRA-full-FT-100k~\cite{longlora}, as well as YaRN-LLaMA and YaRN-Mistral~\cite{yarn}.

\subsection{Main Results}

\label{sec:mainresult}
\textbf{Language modeling over long sequences up to 256k}. Our initial analysis involves benchmarking against advanced long-context LLMs for evaluation lengths up to 256k tokens. To highlight the breadth of our approach, we conduct experiments on two benchmarks: the Proof-pile~\cite{pg19} and PG19~\cite{pile} test sets. We assess perplexity across varying context windows, employing a 256-token sliding window for all measurements. For PG19, our evaluation covers the complete test set comprising 100 documents. In the case of Proof-pile, adhering to the protocol from YaRN~\cite{yarn}, we randomly select 10 instances, each exceeding 128k tokens in length.

Table~\ref{tbl:proof-pile} and Table~\ref{tbl:pg19} present a comparison of perplexity for LLaMA2 and Mistral models extended by various interpolation strategies, evaluated on the Proof-pile and PG19 datasets, respectively. Two primary findings emerge: \textbf{(1)} our extended models exhibit a consistently decreasing perplexity as evaluation sequence length increases from 4k to 256k tokens, indicating improved capacity to exploit longer contexts. \textbf{(2)} Remarkably, even when operating within a context window expanded by a factor of 16—a regime where performance at shorter lengths often degrades—our LongRoPE-2048k models surpass leading baselines on sequences up to 256k tokens in length.

\textbf{Long sequence language modeling beyond 2000k}. To assess performance on exceptionally lengthy documents, we employ the Books3~\cite{pile} dataset. For efficient evaluation, 20 books are randomly chosen, each surpassing 2048k tokens in length, utilizing a sliding window of 256k for processing.

As indicated in Table~\ref{tbl:books3}, LongRoPE is capable of broadening the context windows of LLaMA2-7B and Mistral-7B models up to 2048k, while maintaining perplexity metrics that match or surpass those of the baseline models across shorter sequences ranging from 8k to 128k. Distinct performance patterns emerge between the 2048k variants of LLaMA2 and Mistral. Specifically, Mistral exhibits superior results at limited context lengths but its perplexity rises above 7 once the context exceeds 256k tokens. The results for LLaMA2 are consistent with anticipated outcomes: as the context length grows, perplexity largely declines, aside from slight upticks at the 1024k and 2048k marks. Additionally, with LLaMA2, LongRoPE-2048k yields improved performance when the fine-tuning is conducted at a window length of 256k instead of 128k, attributable to the lower secondary extension ratio employed (8$\times$ versus 16$\times$). Conversely, Mistral demonstrates enhanced outcomes when using a 128k fine-tuning window. This can primarily be explained by the fact that, for Mistral’s 128k and 256k fine-tuning scenarios, we adopt YaRN’s protocol of utilizing a 16k training window, consequently limiting Mistral’s capacity to further enlarge its context window post fine-tuning.

\textbf{Passkey retrieval}. In this section, we analyze the model's usable context window during generation by employing a synthetic passkey retrieval task introduced by~\cite{passkey}. In this evaluation, the model must extract a randomly generated five-digit number (the passkey) that has been embedded at a random position within an extended document. The prompt template used for this task is provided in the appendix. We conduct 10 rounds of this passkey retrieval experiment, ensuring that the passkey’s position is selected uniformly at random within the full evaluation context range for each trial.

As illustrated in Fig.~\ref{fig:passkey}, retrieval accuracy for baseline models sharply declines to zero when the sequence length exceeds 128k tokens. Conversely, our LongRoPE-LLaMA2-2048k (ft=256k) demonstrates robust performance, achieving retrieval accuracies of at least 90\% across a wide range—from 4k up to 2048k tokens—even when extracting a passkey amid millions of tokens. For LongRoPE-Mistral-2048k (ft=128k), perfect retrieval (100\%) is sustained up to 1800k tokens before falling to 60\% accuracy at 2048k, a trend that is consistent with the increased perplexity observed at 2048k in Table~\ref{tbl:books3}.

\textbf{Standard benchmarks within original context window}. The performance of LongRoPE-2048k models is assessed on their native context window utilizing the Hugging Face Open LLM Leaderboard~\cite{huggingfaceboard} in both zero-shot and few-shot evaluation paradigms. Specifically, we adopt a 25-shot setting for ARC-Challenge~\cite{allenai:arc}, a 10-shot configuration for HellaSwag~\cite{zellers2019hellaswag}, a 5-shot approach for MMLU~\cite{mmlu}, and a 0-shot regime for TruthfulQA~\cite{lin2021truthfulqa}.

Table~\ref{tbl:commonsense} demonstrates that our models perform on par with the original benchmark, which was intended for use with shorter context windows, and surpass the original Mistral on TruthfulQA by a margin of +0.5\%. Although LongRoPE-LLaMA2-2048k—fine-tuned on 256k tokens—exhibits a modest drop in performance, it still maintains competitive results across the majority of evaluated tasks.

\subsection{Ablation Results}

\textbf{Effectiveness of the second positional interpolation}. Our progressive extension framework incorporates a search-based approach to apply a secondary, non-uniform positional interpolation to LLMs that have already undergone fine-tuning and extension. To assess its utility, we experiment with our LLaMA2-256k model that has been fine-tuned. The model is further extended to 512k, 1024k, and 2048k token contexts using PI and YaRN methods. As indicated in Table~\ref{tbl:secondsearch}, the non-uniform positional interpolation approach maintains perplexity at a stable level. By comparison, both PI and YaRN methods result in rapidly escalating perplexity as the extension ratio increases.

\textbf{Recovery at Reduced Context Lengths.} To address declines in performance at shorter context windows, we recalibrate the RoPE factors in LongRoPE-2048k by utilizing our search procedure. In this approach, we restrict the upper limit of scale factors during the search, which promotes reduced interpolation for short sequence lengths of 4k and 8k. Table~\ref{tbl:shortperformance} presents the perplexity scores for LongRoPE-LLaMA2-2048k evaluated on Proof-pile at 4k and 8k, in addition to the mean LLM benchmark accuracy. The data indicate that this adjustment yields considerable gains in performance for short contexts.

\textbf{Analysis on the two forms of non-uniformities}. To further understand the contribution of each type of non-uniformity, we conducted ablation studies targeting these components individually. Two sets of experiments were designed: (i) LLaMA2-7B was extended to context lengths of 16k and 32k using multiple strategies—PI, optimization over the RoPE dimension alone, and simultaneous search over both non-uniformities; (ii) a variant of LLaMA2 fine-tuned on 256k context was expanded to 2048k by employing analogous methods. We assessed perplexity for all scenarios without applying additional fine-tuning. According to the results in Table~\ref{tbl:searchdimension}, introducing non-uniformity in the RoPE dimension leads to marked perplexity reductions compared to the conventional linear interpolation of PI. Incorporating non-uniformity into token position substantially enhances results at 16k and 32k, though its effect diminishes at 2048k context length, potentially due to the unprecedented scale. In this extremely long-range regime, simply retaining original tokens without performing interpolation ceases to yield benefit—addressing this limitation is proposed as future research.

\section{Related Works}

Beyond position interpolation techniques, this section reviews alternative strategies explored in related research.

\textbf{Retrieval-based} methods employ an external memory system to retain information from distant contexts, alongside retrieval mechanisms to fetch pertinent documents during inference~\cite{longllama,wang2023augmenting,borgeaud2022improving}. Such systems often require deliberate and explicit changes to the architectures of LLMs. By contrast, our method is more efficient, relying on only minimal adjustments to positional embeddings. Furthermore, our approach accommodates a broader range of long-context applications beyond simple retrieval, including tasks like extensive document summarization and few-shot learning.

\textbf{Attention-based context window extensions}. In addition to positional embedding interpolation, other studies have extended input context within the initial LLM context window by modifying the attention mechanism~\cite{han2023lminfinite, streamingllm,parallelwindow}. The primary strategy involves alleviating the surge in attention demands resulting from new positions through the introduction of innovative attention masks. These strategies are not mutually exclusive, and can work in tandem with positional interpolation approaches.

\textbf{Fine-tuning based approaches} explore strategies for adapting pre-trained LLMs by revising position embeddings, thereby extending their effective context length. Examples include Code LLaMA~\cite{codellama}, LLaMA2 Long~\cite{xiong2023effective}, and ScaledRoPE~\cite{liu2023scaling}, which select a substantially increased base parameter for RoPE and subsequently fine-tune the model for the new target context size. In contrast, our method provides adaptable support for diverse target lengths and is capable of surpassing a 2M token limit. Given that extending fine-tuning to longer contexts (exceeding 128k tokens) is highly GPU-intensive, more recent approaches such as LongLoRA~\cite{longlora} and PoSE~\cite{zhu2023pose} have been introduced to reduce computational costs. Our technique operates independently and is compatible with these efficient fine-tuning methods.

\section{Conclusion}

This paper introduces LongRoPE, a technique that dramatically increases the context length of LLMs to a record-breaking 2048k, all while preserving their performance within the standard, shorter context range. Our approach leverages two distinct types of non-uniformities present in RoPE positional embeddings, identified through an efficient evolutionary search process. This methodology confers dual advantages: it supplies optimal initialization for subsequent fine-tuning and supports an 8$\times$ expansion of the context window even without fine-tuning. Furthermore, we outline a stepwise extension method that utilizes LLMs fine-tuned with a 256k-length context to achieve a 2048k window size, without the necessity of further fine-tuning. Comprehensive experiments corroborate the efficacy of LongRoPE. We anticipate that the LongRoPE-2048k models will unlock a wide array of new applications requiring long-context capabilities and encourage additional advancements in the field.

\section*{Broader Impacts} The research outlined in this paper aims to contribute to progress in Machine Learning as a discipline. While our work may have a variety of societal ramifications, we do not identify any that require particular emphasis at this time.

	\balance

\end{document}